\documentclass[10pt,a4paper]{amsart}
\usepackage[margin=19mm]{geometry}
\usepackage{amsmath,amssymb,mathtools}
\usepackage[scr=rsfs,cal=euler]{mathalfa}
\usepackage{xfrac}
\usepackage[unicode,hidelinks,bookmarks=false]{hyperref}
\newcommand{\ie}{i.e.\ }
\newcommand{\cf}{cf.\ }
\newcommand{\resp}{respectively\ }
\newcommand{\Subsection}{\subsection}
\providecommand{\nobreakdash}{\nobreak\hbox{-}}
\setlength{\emergencystretch}{2em}
\multlinegap0pt
\usepackage{graphicx}
\numberwithin{equation}{section}
\renewcommand*\P{\mathcal{P}}
\newcommand*\M{\mathcal{M}}
\newcommand*\cb{\mathcal{B}}
\newcommand*\cq{\mathcal{Q}}
\newcommand*\cF{\mathcal{F}}
\newcommand*\cT{\mathcal{T}}
\newcommand*\cC{\mathcal{C}}
\newcommand*\cg{\mathcal{G}}
\newcommand*\cl{\mathcal{L}}
\newcommand*\cm{\mathcal{M}}
\newcommand*\cN{\mathcal{N}}
\newcommand*\cR{\mathcal{R}}
\newcommand*\cU{\mathcal{U}}
\newcommand*\cV{\mathcal{V}}
\newcommand*\cw{\mathcal{W}}
\newcommand*\cv{\mathcal{V}}
\newcommand*\cz{\mathcal{Z}}
\newcommand*\cO{\mathcal{O}}
\newcommand*\cal{\mathcal}

\newcommand{\bfone}{\mathbf{1}}
\newcommand{\rId}{\mathrm{Id}}
\newcommand{\DMI}{\mathrm{DMI}}
\newcommand{\MI}{\mathrm{MI}}
\newcommand{\rCard}{\mathrm{Card}}
\newcommand{\rmD}{\mathrm{D}}
\newcommand{\Graph}{\mathrm{graph}}
\newcommand{\rdet}{\mathrm{det}}
\newcommand*\R{\mathbb{R}}
\newcommand*\Z{\mathbb{Z}}
\newcommand*\N{\mathbb{N}}
\newcommand*\C{\mathbb{C}}
\newcommand*\T{\mathbb{T}}
\newcommand{\bbP}{\mathbb{P}}
\renewcommand*\dim{{\rm dim}\,}
\DeclareMathOperator{\Iindex}{index}
\newcounter{stepcount}
\newenvironment{step}{\refstepcounter{stepcount}\begin{proof}[Step~\thestepcount]}{\let\qed\relax\end{proof}}
\newcommand*{\mk}{\mkern -1mu}
\newcommand*{\Mk}{\mkern -2mu}
\newcommand*{\mK}{\mkern 1mu}
\newcommand*{\MK}{\mkern 2mu}
\newcommand*{\romanenumi}{\renewcommand*{\theenumi}{\roman{enumi}}}
\newcommand*{\Romanenumi}{\renewcommand*{\theenumi}{\Roman{enumi}}}
\newcommand*{\alphenumi}{\renewcommand*{\theenumi}{\alph{enumi}}}
\newcommand*{\Alphenumi}{\renewcommand*{\theenumi}{\Alph{enumi}}}
\let\oldtilde\tilde
\renewcommand*{\tilde}[1]{\mathchoice{\widetilde{#1}}{\widetilde{#1}}{\oldtilde{#1}}{\oldtilde{#1}}}
\let\oldhat\hat
\renewcommand*{\hat}[1]{\mathchoice{\widehat{#1}}{\widehat{#1}}{\oldhat{#1}}{\oldhat{#1}}}
\let\oldforall\forall
\renewcommand*{\forall}{\mathrel{\oldforall}}
\newtheorem{theo}{Theorem}[section]
\newtheorem{prop}[theo]{Proposition}
\newtheorem{lemm}[theo]{Lemma}
\newtheorem{coro}[theo]{Corollary}
\theoremstyle{definition}
\newtheorem{defi}[theo]{Definition}
\newtheorem{nota}[theo]{Notation}
\newtheorem{claim}[theo]{Claim}
\theoremstyle{remark}
\newtheorem{rema}[theo]{Remark}
\newenvironment{enonce}[1]{\begin{claim}}{\end{claim}}
\title[Generating functions: Maslov index equals fiber-Hessian Morse-index difference]{Maslov index of a Lagrangian arc with a generating function and vertically transverse endpoints equals its fiber-Hessian Morse-index difference}
\author{Marie-Claude Arnaud, Anna Florio and Valentine Roos}
\date{Source: Annales Henri Lebesgue \textbf{7} (2024), 307--355; DOI 10.5802/ahl.202; official published version}
\begin{document}
\maketitle
\noindent\textit{Editorial scope.} The counted unit is original Theorem 3.3 with its complete authored proof, which contains the complete proofs of Lemma 3.4 and Lemma 3.5. Retained uncounted core: the whole of the paper's own Section 2 "On Maslov index" (the definitions of the cooriented singular cycle, of the height of one Lagrangian above another, of the Maslov index of an arc, and the invariance of the Maslov index under symplectic reduction, with the complete proofs of Propositions 2.3, 2.13, 2.14, 2.15, 2.17 and of Lemmas 2.19 and 2.21), plus Definition 3.1 of a generating function and Notation 3.2. Everything is copied exactly from the publisher-native source. The introduction, the dynamical and angular Maslov index material, the graph-selector and twist material of Sections 4--5 and all applications are not selected. The published theorem, definition, notation and equation numbers are reproduced.\par
\noindent\textit{Wrapper note.} The publisher class supplies the unnumbered wrapper \texttt{enonce}; in the retained core it is used only as \texttt{enonce\{Claim\}}, so it is delivered as the numbered Claim environment that prints the published "Claim 2.16/2.18/2.20". No mathematical content, hypothesis, constant or conclusion is altered.
\setcounter{section}{1}
\section{On Maslov index}\label{section def MI}

\subsection{Some reminders on Maslov index}\label{ssremindersMslov}

Let $\cal M$ be a $2d$-dimensional symplectic manifold that admits a Lagrangian foliation $\cal V$. We denote by $V(x)=V_x:=T_x\cal V$ its associated Lagrangian bundle. Let $p:T\M\rightarrow \cal M$ be the canonical projection. Let $\Lambda(\cal M)$ be the Grassmannian of Lagrangian subspaces of $T\cal M$. We recall that $\Lambda(\cal M)$ is a smooth manifold with dimension~$2d+\frac{d(d+1)}{2}$. The fibered singular cycle associated to $\cv$ is the set
\[
\Sigma(\cal M)=\left\{ L\in \Lambda(\cal M) :\ L\cap V_{p(L)}\not=\{ 0\}\right\}.
\]
Every fiber $\Sigma_x(\cal M)$ of $\Sigma(\cal M)$ is a cooriented algebraic singular hypersurface of $\Lambda_x(\cal M)$, see e.g.~\cite{McDuffSalamon2017,RobbinSalamon1993}. Hence $\Sigma(\cal M)$ is a cooriented singular hypersurface of $\Lambda(\cal M)$.

The singular locus of $\Sigma(\cal M)$ is then $\{ L\in\Lambda(\cal M) :\ \dim(L\cap V_{p(L)})\geq 2 \}$ and the regular locus is
\begin{equation}\label{Esinglocus}
\Sigma_1:=\left\{ L\in\Lambda :\ \dim\left(L\cap V_{p(L)}\right)=1 \right\}.
\end{equation}
Once a coorientation of $\Sigma(\cal M)$ is fixed, it is classical to associate to every continuous loop $\Gamma:\T\rightarrow \Lambda (\cal M)$ its \emph{Maslov index} $\MI(\Gamma)$, that satisfies the following properties:
\begin{itemize}
\item two homotopic loops have the same Maslov index;
\item if $\Gamma$ is a loop that avoids the singular locus and is topologically transverse to the regular one of $\Sigma(\cal M)$, then $\MI(\Gamma)$ is the number of signed intersections of $\Gamma$ with $\Sigma(\cal M)$ with respect to the chosen coorientation;
\item every loop is homotopic to a smooth loop that avoid the singular locus and is transverse to the regular locus.
\end{itemize}
An arc in $\Lambda(\cal M)$ is an immersion $\Gamma:[0,1]\to\Lambda(\cal M)$. In particular, $\Gamma([0,1])$ does not have self-intersections. By smooth arc, we mean a $C^\infty$ arc. When $\Gamma: [0, 1]\rightarrow \Lambda(\cal M)$ is an arc whose endpoints are in $\Lambda(\M)\backslash \Sigma(\M)$, following Duistermaat~\cite[p.~183]{Dui76}, we can concatenate $\Gamma$ with an arc $\Gamma_1$ that connects $\Gamma (1)$ to $\Gamma(0)$ in $\Lambda(\M)\backslash \Sigma(\M)$. The Maslov index of $\Gamma$ is the Maslov index of this loop, which is independent from the choice of $\Gamma_1$ { since $\Gamma_1$ is in $\Lambda(\M)\setminus\Sigma(\M)$}.


\begin{rema}\label{rmq path transverse}
If $\Gamma:[0,1]\to\Lambda(\cm)$ is an arc contained in $\Lambda(\cm)\setminus\Sigma(\cm)$, i.e., $\Gamma(t)\cap\Sigma_{p\circ\Gamma(t)}=\{0\}$ for every $t$, then its Maslov index $\MI(\Gamma)$ is zero.
\end{rema}

\subsection{Coorientation of \texorpdfstring{$\Sigma^1$}{Sigma1}}\label{subsec defi MI}


We now give some details concerning the singular and regular loci of $\Sigma (\cal M)$ and explain our choice of coorientation of $\Sigma(\cal M)$. For more details, see for example~\cite{Dui76}. For ease of reading, we denote $\Sigma (\cal M)$ (resp. $\Lambda(\cal M)$) by $\Sigma$ (resp. $\Lambda$).

Then $\Sigma$ is an algebraic subvariety of $\Lambda$ that is the union of
\begin{itemize}
\item the regular locus that is the smooth submanifold of codimension~1 and is defined in~\eqref{Esinglocus},
\item the boundary of $\Sigma_1$, i.e. the singular locus $\Sigma\setminus\Sigma^1$, that is a finite union of submanifolds with codimension at least 3.
\end{itemize}
Since every loop is homotopic to a smooth loop avoiding the singular locus and intersecting transversally the regular one and since two homotopic loops have the same Maslov index, we just have to define the coorientation at points of $\Sigma^1$. To do that, we introduce the notion of height in a symplectic vector space $(E^{2d}, \Omega)$.

We fix a reference Lagrangian subspace $V$ of $E$ and denote by $P^V$ the canonical projection on the quotient vector space $E/V$. If $L_1$, $L_2$ are two Lagrangian subspaces of $E$ that are transverse to $V$, we define the height of $L_1$ above $L_2$ with respect to $V$, see~\cite{Arn08}, as follows.

\begin{defi}
Let $L_1,L_2\subset E$ be two Lagrangian subspaces both transverse to $V$. The \emph{height} of $L_2$ above $L_1$ with respect to $V$ is the quadratic form
\[
\cq_V(L_1,L_2):E/V\to\R
\]
defined by
\[
\forall v\in E/V, \cq_V(L_1,L_2)(v):=\Omega\left(\left({P}^V\middle|_{L_1}\right)^{-1}(v),\left({P}^V\middle|_{L_2}\right)^{-1}(v)\right).
\]
\end{defi}

With the hypotheses of this definition, the kernel of $ \cq_V(L_1,L_2)$ is isomorphic to $L_1\cap L_2$. In particular, $L_1$ is transverse to $L_2$ if and only if $Q_V(L_1,L_2)$ is non degenerate.\newline
We have
\begin{itemize}
\item If $L_1,L_2,L_3$ are Lagrangian subspaces in $E$, all transverse to $V$, it holds, see~\cite{Arn08},
\begin{equation}\label{prop:quadr:form}
\cq_V(L_1,L_3)=\cq_V(L_1,L_2)+\cq_V(L_2,L_3).
\end{equation}
\item if $V$, $K$, $L$ are Lagrangian subspace of $E$ such that each of them is transverse to the two others, then { $\cq_V(K,L)\circ P^V|_{L}=-\cq_K(V,L)\circ P^K|_{L}$ and then} $\cq_V(K,L)$ and $-\cq_K(V,L)$ have the same signature.\newline
Let us prove that $\cq_V(K,L)\circ P^V|_{L}=-\cq_K(V,L)\circ P^K|_{L}$. For $\ell\in L$, there exists a unique pair of vectors $v\in V, k\in K$ such that $\ell=v+k$. then we have
\begin{itemize}
\item $\cq_V(K,L)\circ P^V(\ell)=\Omega(k, \ell)=\Omega(k, v)$;
\item $-\cq_K(V,L)\circ P^K(\ell)=-\Omega(v, \ell)=-\Omega(v, k)=\Omega(k, v) $.
\end{itemize}

\item if $L$ and $K$ are Lagrangian subspaces that are transverse to $V$ and if $\phi:E\,\rotatebox[origin=c]{90}{$\circlearrowleft$}\,$ is a symplectic isomorphism, then $\cq_V(K, L)$ has same signature as $\cq_{\phi(V)}(\phi(K), \phi(L))$.
\end{itemize}
We now describe the local coorientation of $\Sigma_1$ that we will use. Let us fix $L_0\in\Sigma_1$ and let $x_0:=p(L_0)$. We have $\dim(L_0\cap V_{x_0})=1$. We fix a Darboux chart $F=(F_1, F_2): \cU\rightarrow\R^{2d}$ at $x_0$ such that $\cal U$ is a small neighborhood of $x_0$ in $\M$, $F(\cal U)=[a, b]^{d}\times [a, b]^d$ and $DF_{2|L_0}$ is injective and
\[
\forall x\in \cU, F({ \cV(x)}\cap \cal U)= F_1(x)\times [a, b]^d.
\]
Let us explain why such a chart exists. Using~\cite[Theorem~7.1]{Weinstein1971}, we can map locally the foliation $\cV$ onto the vertical foliation of $\R^{2d}$ by a symplectic chart $(U, \Phi)$. Then, composing with a symplectic isomorphism $\psi_t(x, y)=(x, y+tx)$ of $\R^d\times\R^d$, for some $t\in\R$, we obtain a new chart $F=(F_1, F_2)$ that maps $\cV$ onto the vertical foliation such that $DF(L_0)$ is transverse to $\{0\}\times \R^d$ and then $DF_{2|L_0}$ is injective.


We denote by ${\cal K}$ the Lagrangian foliation with leaves $F^{-1}([a, b]^d\times \{y_0\}$). Then it is transverse to the vertical bundle $V$. Moreover, $T_{x_0}{\cal K}$ and $L_0$ are transverse, since ${DF_2}_{\vert L_0}$ is injective. We denote by $K$ the tangent bundle to ${\cal K}$. { Because $\dim(L_0\cap V)=1$,} the kernel of $\cq_{K}(V, L_0)$ is 1-dimensional. We denote by $n$ the index\footnote{The \emph{index} of a quadratic form is the maximum dimension of a subspace of $E$ on which the quadratic form is negative definite.} of $\cq_{K}(V, L_0)$. We define
\begin{align*}
\P_1 &= \left\{ L\in \Lambda\backslash \Sigma;\ p(L)\in \cU,\ L\pitchfork K,\ \Iindex\cq_{K}(V, L)=n\right\}
\\
\intertext{and}
\P_2 &= \left\{ L\in \Lambda\backslash \Sigma;\ p(L)\in \cU,\ L\pitchfork K,\ \Iindex\cq_{K}(V, L)=n+1\right\}.
\end{align*}
Observe that $\P_1$ and $\P_2$ are connected and that $\P_1\cup{\P_2}\cup \Sigma_1$ is a neighbourhood of $L_0$ in $\Lambda$. Hence $\P_1$ and $\P_2$ define locally a coorientation\footnote{Let $\gamma\subset \P_1\cup\P_1\cup \Sigma_1$ be a path from $\P_2$ to $\P_1$, crossing $\Sigma_1$ transversally once at $\gamma(t)$. Then $\gamma'(t)\in \R_+ N$, where $N$ is a normal vector field to $\Sigma_1$ and determines a coorientation of $\Sigma_1$ at $L_0$.} of $\Sigma_1$ at $L_0$. To be sure that we obtain a global coorientation of $\Sigma$, we have to prove that this local coorientation is independent from the choice of our foliation ${\cal K}$. We just have to look at what happens in the fiber $\Lambda_{x_0}$ for different choices of $K_{x_0}$. In other words, we will prove a result in a fixed symplectic vector space $(E, \Omega)$.

\begin{prop}\label{Pcoorientation}
Let $V$, $L_0$ be two Lagrangian subspaces of $E$ such that $\dim (L_0\cap V)=1$. Let $K_1$, $K_2$ be two Lagrangian subspaces of $E$ that are transverse to $L_0$ and $V$. We denote by $n_i$ the index of $\cq_{K_i}(V, L_0)$. There exists a neighbourhood $\cal U$ of $L_0$ in the Lagrangian Grassmannian of E such that
\begin{multline*}
\left\{ L\in \cU;\ L\pitchfork K_1,\ \Iindex\cq_{K_1}(V, L)=n_1+1\right\}=\\
\left\{ L\in \cU;\ L\pitchfork K_2,\ \Iindex\cq_{K_2}(V, L)=n_2+1\right\}
\end{multline*}
and
\begin{multline*}
\left\{ L\in \cU;\ L\pitchfork K_1,\ \Iindex\cq_{K_1}(V, L)=n_1\right\}=\\
\left\{ L\in \cU;\ L\pitchfork K_2,\ \Iindex\cq_{K_2}(V, L)=n_2\right\}.
\end{multline*}
\end{prop}

\begin{proof}
Because, for $j=1,2$, $V$ and $K_j$ are transverse Lagrangian subspa\-ces, every basis $(e_1, \,\dots,\, e_d)$ of $V$ can be completed in a symplectic basis $(e, f^j)=(e_1, \,\dots,\, e_d, f^j_1,\, \dots,\, f^j_d)$ of $E$ such that $f^j_i\in K_j$. Then, every Lagrangian subspace $L$ of $E$ that is close enough to $L_0$ is the graph in this basis of a $d\times d$ symmetric matrix $S_j^L$ that continuously depends on $L$ and is close to $S_j^{L_0}$.  We identify $V$ with $\R^d$ via the basis $(e_i)$.\newline
As $\dim\ker S_j^{L_0}=1$, we have $\R^d=\R \ell_j(L_0)+E_j(L_0)$ where $\ker S_j^{L_0}=\R\ell_j(L_0)$ and $E_j(L_0)=(\R \ell_j(L_0))^\bot$ is the orthogonal of $\ker S_j^{L_0}$ for the usual euclidean scalar product, i.e., the sum of the eigenspaces for the non-zero eigenvalues. Observe that we can choose $\ell_1(L_0)=\ell_2(L_0)$. For $L$ in some neighbourhood $\cU$ of $L_0$, $S_j^L$ has a spectral gap with one eigenvalue $\lambda(S_j^L)$ close to $0$ and the others far away from $0$. Hence we can continuously extend $\ell_j(L)$ and $E_j(L)$ for $L$ close to $L_0$ in such a way that $\ell_j(L)$ is an eigenvector for the eigenvalue that is close to $0$, and $E_j(L)$ is $(\R\ell_j(L))^\perp$. Moreover, the signature of the restriction of $S_j^L$ to $E_j(L)$ remains equal to its value for $L=L_0$ if $\cal U$ is small enough. \newline
The matrix of $\cq_{K_j}(V, L)$ in the basis $(P^{K_j}(e_1), \,\dots,\, P^{K_j}(e_d))$ of $E/K_j$ is $S_j^L$ and then to estimate the index of $\cq_{K_j}(V, L)$, we only need to know the sign of $\lambda(S_j^L)$.
We recall that when $L\in \cU$ is transverse to $V$, we have $\cq_{K_j}(V, L)\circ (P^{K_j}|_{L})^{-1}=-\cq_{V}(K_j, L)\circ(P^V|_{L})^{-1}$. The matrix of $-\cq_{V}(K_j, L)$ in the basis $(P^V(f_1^j),\,\dots,\linebreak P^V(f_d^j))$ is $(S_j^L)^{-1}$ and thus we are reduced to estimate the sign of the eigenvalue of $(S_j^L)^{-1}$ that has the largest absolute value. Observe that $P^V(f_i^1)=P^V(f_i^2)$. We denote by $S$ the matrix of $ \cq_{V}(K_1,K_2)$ in the same basis and we deduce from~\eqref{prop:quadr:form} that
\begin{equation}\label{Ematsym}
-\Big(S_1^L\Big)^{-1}=-\Big(S_2^L\Big)^{-1}+S
\end{equation}
Let us denote by $\|\cdot\|_2$ the usual Euclidean norm on $\R^d$ and let us endow the set of $d$-dimensional matrices with the associated norm defined by
\[
\| S\|=\sup_{\|v\|_2=1} \|Sv\|_2\,.
\]
Then if $\,\cU$ is small enough, there exists $C>\| S\|$ such that for every $L\in \cU$,
$(\lambda(S_j^L))^{-1}$ is the only eigenvalue of $(S_j^L)^{-1}$ whose absolute value is larger than $3C$ and $C$ is an upper bound of the modulus of all the other eigenvalues of $(S_j^L)^{-1}$.

Let us prove that $\lambda(S_1^L)$ and $\lambda(S_2^L)$ have the same sign. Let $v\in\R^d$ be an eigenvector of $S_1^L$ for the eigenvalue $\lambda(S_1^L)$. Then there exists $v_1, v_2\in\R^d$ that are mutually orthogonal such that {$v=v_1+v_2$,} $S_2^Lv_1=\lambda(S_2^L)v_1$ and $v_2$ is orthogonal to the eigenspace of $S_2^L$ for $\lambda(S_2^L)$. Using~\eqref{Ematsym}, we obtain
\[
\big(\lambda\left(S_1^L\right)\big)^{-1}-\big(\lambda\left(S_2^L\right)\big)^{-1}\frac{\| v_1\|_2^2}{\| v\|_2^2}=\frac{v_2^T}{\| v\|_2} \big(S_2^L\big)^{-1}\frac{v_2}{\| v\|_2}- \frac{v^T}{\| v\|_2}S\frac{v}{\| v\|_2}
\]
Observe that the absolute value of the right-hand term is less than $2C$. If $\lambda(S_1^L)$ and $\lambda(S_2^L)$ have different signs, then the absolute value of the left-hand term is larger than the absolute value of $(\lambda(S_1^L))^{-1}$, then larger than $3C$, which provides a contradiction.
\end{proof}

In order to define the Maslov index, we first introduce the notion of positive (resp. negative) arcs. Recall that $K$ denotes the tangent bundle to $\cal K$, where $\cal K$ is the Lagrangian foliation with leaves $F^{-1}([a,b]^d\times\{y_0\})$.

\begin{defi}
With the same notation, an arc $\Gamma:(-\varepsilon_0, \varepsilon_0)\rightarrow \Lambda$ such that
\[
\Gamma((-\varepsilon_0,\varepsilon_0))\cap \Sigma=\Gamma((-\varepsilon_0,\varepsilon_0))\cap \Sigma_1=\{ \Gamma(0)\} =\{L_0\}
\]
and that is topologically transverse to $\Sigma^1$ is {\sl positive} if there exists $\varepsilon>0$ such that
\begin{itemize}
\item for every $t\in (-\varepsilon, 0)$, ${\rm index}(\cq_{K}(V, \Gamma(t)))={\rm index}(\cq_{K}(V, L_0))+1$;
\item for every $t\in (0, \varepsilon)$, ${\rm index}(\cq_{K}(V, \Gamma(t)))={\rm index}(\cq_{K}(V, L_0))$.
\end{itemize}
Respectively, an arc $\Gamma:(-\varepsilon_0,\varepsilon_0)\to\Lambda$ is {\sl negative} if $\Gamma\circ(-\rId)$ is positive.
\end{defi}

\begin{rema}
This is equivalent to
\begin{itemize}
\item for every $t\in (-\varepsilon, 0)$, ${\rm index}(\cq_{V}(K, \Gamma(t)))=d-{\rm index}(\cq_{K}(V, L_0))-1$;
\item for every $t\in (0, \varepsilon)$, ${\rm index}(\cq_{V}(K, \Gamma(t)))=d-{\rm index}(\cq_{K}(V, L_0))$.
\end{itemize}
\end{rema}

\begin{defi}
Let $\Gamma:[a, b]\rightarrow \Lambda$ be an arc.
\begin{itemize}
\item A $t\in[a,b]$ is a {\sl crossing} for $\Gamma$ if $\Gamma(t)\in\Sigma$.
\item The arc $ \Gamma$ is in {\sl general position} with respect to $\Sigma$ if $\Gamma(a),\Gamma(b)\in\Lambda\setminus\Sigma$ and the path $\Gamma$ is { topologically} transverse to $\Sigma$. {
\item The arc $ \Gamma$ is in {\sl D-general position} with respect to $\Sigma$ if $\Gamma(a),\Gamma(b)\in\Lambda\setminus\Sigma$ and the path $\Gamma$ is transverse (in the differentiable sense) to $\Sigma$.}
\end{itemize}
\end{defi}

\begin{rema}
If $\Gamma:[a,b]\to\Lambda$ is in general position with respect to $\Sigma$, then each crossing for $\Gamma$ is isolated. {Let $[a, b]$ be fixed and let $k\in\N^*\cup\{ \infty\}$. Then, the set of $C^k$ arcs $\Gamma:[a, b]\to \Lambda$ that are in D-general position with respect to $\Sigma$ is open for the $C^1$-topology.}
\end{rema}

Let $\Gamma:[a,b]\to\Lambda$ be an arc in general position with respect to $\Sigma$. A crossing $t$ is called \emph{positive}, respectively \emph{negative}, if there exists $\epsilon>0$ such that the arc $\Gamma\vert_{[t-\epsilon, t+\epsilon]}:[t-\epsilon, t+\epsilon]\to\Lambda$ is positive, respectively negative.



\begin{defi}\label{defi MI for arc}
Let $\Gamma:[a,b]\to\Lambda$ be an arc in general position with respect to $\Sigma$. The Maslov index of $\Gamma$ with respect to {$V$ or $\cV$} is
\begin{multline*}
\MI(\Gamma) :=
\\
\rCard\{t :\, t\text{ is a positive crossing for }\Gamma\}-\rCard\{t:\, t \text{ is a negative crossing for }\Gamma\}\,.
\end{multline*}
\end{defi}

The notion of Maslov index can be extended to Lagrangian paths that are not in general position.
\begin{defi}
Let $\Gamma:[a,b]\to\Lambda$ be a path such that $\Gamma(a),\Gamma(b)\in\Lambda\setminus\Sigma$ (not necessarily in general position with respect to $\Sigma$). Let { $\tilde\Gamma:[a,b]\to \Lambda$ be a smooth arc that is $C^1$-close to $\Gamma$ and is in general position with respect to $\Sigma$. } Then
\[
\MI(\Gamma):=\MI(\tilde\Gamma)\,.
\]
\end{defi}

For the existence of the perturbation $\tilde\Gamma$ of $\Gamma$ and for the independence of the previous definition from the choice of $\tilde\Gamma$ we refer to~\cite{Arn72} or~\cite{Capleem}.


\begin{rema}\label{rmk invariance diffeo cs}
Let $\phi$ be a conformally symplectic diffeomorphism on $\M$. Let $\Gamma:[a.b]\to\Lambda$ be a smooth path such that $\Gamma(a),\Gamma(b)\notin\Sigma$. Then
\[
D\phi(\Gamma):[a,b]\ni t\to D\phi(\Gamma(t))\in\Lambda
\]
is still a smooth path such that $D\phi(\Gamma)(a),D\phi(\Gamma)(b)$ do not belong to
\[
\left\{ L\in\Lambda :\ L\cap D\phi(V)_{p(L)}\neq\{0\}\right\}\,,
\]
where $D\phi(V)_x$ is the tangent bundle associated to the Lagrangian foliation $\phi(\cv)$. Then the Maslov index $\MI(\Gamma)$, calculated with respect to the Lagrangian foliation $\cv$, is equal to the Maslov index $\MI(D\phi(\Gamma))$, calculated with respect to the Lagrangian foliation $\phi(\cv)$.

Moreover, if $\phi(\cv)=\cv$, then
\[
\MI(\Gamma)=\MI(D\phi(\Gamma))\,.
\]
In particular, for $\M=T^*M$, the Maslov index is invariant under \emph{vertical translations}, that is by any diffeomorphism of the form $\phi(p)=p+\eta\circ\pi(p)$, where $\eta$ is a closed 1-form on $T^*M$.
\end{rema}

\subsection{Dynamical Maslov index}

We now give the definition of dynamical Maslov index.

\begin{defi}
Let $(\M,\omega)$ be a symplectic manifold that admits a Lagrangian foliation $\cv$. Let $(\phi_t) $ be an isotopy of conformally symplectic diffeomorphisms of $\M$. Let $L\in\Lambda$ and $[\alpha,\beta]\subset \R$ be such that $D\phi_\alpha(L),D\phi_\beta(L)\notin\Sigma$. Then
\[
\DMI\left(L,(\phi_t)_{t\,\in\,[\alpha,\beta]}\right):= \MI(\Gamma)\,,
\]
where $\Gamma$ is the Lagrangian path $[\alpha,\beta]\ni t\mapsto \Gamma(t):=D\phi_t(L)\in\Lambda$ and the Maslov index $\MI(\Gamma)$ is calculated with respect to the Lagrangian foliation $\cv$.
\end{defi}



\subsection{Twist and Maslov index}


In this section, we work in $T^*M$ and we denote ${\cal V}(x)=T_x^*M$.

In the introduction, we gave the definition of an isotopy which twists the vertical. We can enhance this in the following way (we adopt the same notations $F$, $g_t$ and $\cg_t$ as in the definition of twist of the vertical).

\begin{defi}
An isotopy $(\phi_t)$ of conformally symplectic diffeomorphisms of $\M$ \emph{strictly} twists the vertical if it twists the vertical and at every $t_0\in\R$
\begin{itemize}
\item for all $t\in(t_0,t_0+\epsilon)$ the image $F(\cg_t)$ is the graph of a function $p\mapsto q=dg_t(p)$ where $g_t$ is a strictly convex function, i.e., such that $d^2g_t$ is positive definite;
\item for all $t\in(t_0-\epsilon,t_0)$ the image $F(\cg_t)$ is the graph of a function $p\mapsto q=dg_t(p)$ where $g_t$ is a strictly concave function, i.e., such that $d^2g_t$ is negative definite.
\end{itemize}
\end{defi}




Observe that the condition of convexity depends on the charts we choose (even if the property of twisting the vertical is invariant under symplectic conjugation that preserves the vertical foliation). This is a motivation to provide a result on the twist property that doesn't rely on any specific chart.

\begin{prop}\label{Ptwistintrinsec}
Let $(\phi_t)$ be an isotopy of conformally symplectic diffeomorphisms of $T^*M$ that twists the vertical. Let $x\in T^*M$ and let $t_0\in\R$. We denote $x_t=\phi_t(x)$. Let $K$ be a continuous Lagrangian bundle that is defined in a neighbourhood of $x_{t_0}$ and is transverse to the vertical bundle. Then, there exists $\varepsilon>0$ such that
\begin{itemize}
\item $\forall t\in (t_0, t_0+\varepsilon), Q_{K(x_t)}(D\phi_t\circ (D\phi_{t_0})^{-1}V(x_{t_0}), V(x_t)) $ is a negative semi-definite quadratic form;
\item $\forall t\in (t_0-\varepsilon, t_0), Q_{K(x_t)}(D\phi_t\circ (D\phi_{t_0})^{-1}V(x_{t_0}), V(x_t)) $ is a positive semi-definite quadratic form.
\end{itemize}
Moreover, when $(\phi_t)$ strictly twists the vertical, the considered quadratic forms are negative definite or positive definite.
\end{prop}


\begin{proof}[Proof of Proposition~\ref{Ptwistintrinsec}]
We fix $\varepsilon>0$ and a vertically foliated chart $F=(F_1, F_2): \cU\rightarrow\R^{2d}$ such that $x_{t_0}\in\cal U$, $F(x_{t_0})=0$ and for $t\in (t_0-\varepsilon, t_0+\varepsilon)$
\begin{itemize}
\item ${\cg}_t:=(\phi_t\circ\phi_{t_0}^{-1})(F^{-1}(\{ 0_{\R^d}\}\times (-\frac{a}{2}, \frac{a}{2})^d)) \subset\cU$;
\item $F({\cg_t})$ is the graph of a function $p\mapsto q=dg_t(p)$ where
\begin{enumerate}
\item for $t\in [t_0, t_0+\varepsilon)$, $g_t$ is a convex function;
\item for $t\in (t_0-\varepsilon, t_0]$, $g_t$ is a concave function.
\end{enumerate}
\end{itemize}
As previously, we denote by ${\cal K}$ the Lagrangian foliation with leaves $F^{-1}([-a, a]^d\times \{y_0\}$) and by $K$ its tangent bundle. For $t\in(t_0,t_0+\epsilon)$ (resp. $(t_0-\epsilon, t_0)$), the quadratic form
\[
\cq_{K(x_t)}\left(D\left(\phi_t\circ \phi_{t_0}^{-1}\right)V(x_{t_0}), V(x_t)\right)=-\cq_{K(x_t)}\left(V(x_t), D\left(\phi_t\circ \phi_{t_0}^{-1}\right)V(x_{t_0})\right)
\]
expressed in the chart $F$ is just $-D^2g_t(F_2(x_{t_0}))$ that is a negative (resp. positive) semi-definite quadratic form because the isotopy twists the vertical.\newline
When the isotopy strictly twists the vertical, we obtain in this case a negative (resp. positive) definite quadratic form.

Observe that the bundle $K$ that we use in the proof is not necessarily the same bundle as in the statement. But because the two are transverse to the vertical foliation and we consider the height between the vertical and Lagrangian subspaces that are close to the vertical ($\varepsilon$ is small), the two indices are the same (we can build an isotopy between the two bundles that won't change the signature).
\end{proof}

\begin{prop}\label{Ptermchpvect}
Let $(X_t)$ be a conformally symplectic vector field that generates an isotopy of conformally symplectic diffeomorphisms of $T^*M$. We assume that for every $x\in T^*M$, there exists a vertically foliated chart $F=(F_1, F_2):\cU\to \R^{2d}$ such that when we express the vector field $X=(X_q, X_p)$ in this chart, then $\partial_pX_q$, which is always symmetric due to the symplectic nature of the vector field, is positive definite.



Let $x\in T^*M$ and denote $x_t=\phi_t(x)$. Let $t_0\in\R$. Let $K$ be a continuous Lagrangian bundle defined in a neighbourhood of $x_{t_0}$ and transverse to the vertical bundle. Then there exists $\varepsilon>0$ such that
\begin{itemize}
\item $\forall t\in (t_0, t_0+\varepsilon), Q_{K(x_t)}(D\phi_t\circ (D\phi_{t_0})^{-1}V(x_{t_0}), V(x_t))$ is positive definite;
\item $\forall t\in (t_0-\varepsilon, t_0), Q_{K(x_t)}(D\phi_t\circ (D\phi_{t_0})^{-1}V(x_{t_0}), V(x_t))$ is negative definite.
\end{itemize}
\end{prop}

\begin{proof}[Proof of Proposition~\ref{Ptermchpvect}]
As noticed in the proof of Proposition~\ref{Ptwistintrinsec}, we only need to prove the result for the tangent space $K$ to Lagrangian foliation ${\cal K}$ with leaves $F^{-1}([-a, a]^d\times \{y_0\}$).

In the chosen chart, the Jacobian matrix of $X$ is
\[
DX(x_{t_0})=
\begin{pmatrix}\partial_qX_q(x_{t_0})&\partial_pX_q(x_{t_0})\\
\partial_qX_p(x_{t_0})&\partial_pX_p(x_{t_0})
\end{pmatrix}
\]
and if we denote
\[
D\phi_t(x) \big(D\phi_{t_0}\big)^{-1}(x_{t_0})=
\begin{pmatrix}a_t&b_t\\ c_t&d_t
\end{pmatrix}
\]
then we have
\begin{equation}\label{Ehamlin}
\begin{cases}
\dot b_t=\partial_qX_qb_t+\partial_pX_qd_t\\
\dot d_t=\partial_qX_pb_t+\partial_pX_pd_t\,.
\end{cases}
\end{equation}
Hence uniformly in $x$ it holds $d_t= \bfone_d+o(t-t_0)$ and $b_t=(t-t_0)\partial_pX_q(x_{t_0})+o((t-t_0)^2)$, which gives $b_t(d_t)^{-1}= (t-t_0)\partial_pX_q(x_{t_0})+o((t-t_0)^2)$. Because $b_t(d_t)^{-1}$ is the matrix of
\[
Q_{K(x_t)}\left(D\phi_t\circ \big(D\phi_{t_0}\big)^{-1}V(x_{t_0}), V(x_t)\right)
\]
in the chart, this gives the wanted result.
\end{proof}



\begin{prop}\label{PtwistMAslov}
Let $(\phi_t)$ be an isotopy of conformally symplectic diffeomorphisms of $T^*M$ that twists the vertical. Then if $L\in\Lambda$ and $[\alpha, \beta]\subset \R$ are such that $D\phi_\alpha(L), D\phi_\beta(L)\notin \Sigma$, then
\[
\DMI\left(L, (\phi_{ t})_{t\,\in\,[\alpha, \beta]}\right)\leq 0.
\]
\end{prop}

\begin{proof}[Proof of Proposition~\ref{PtwistMAslov}]
Let us first assume that $(\phi_t)$ is an isotopy {satisfying the conclusion of Proposition\ref{Ptwistintrinsec} (with definite quadratic forms) in a neighborhood of $(D\phi_tL)_{t\,\in\,[\alpha,\beta]}$.} Perturbing $L$, we can assume that $(D\phi_{ t}L)_{t\,\in\,[\alpha, \beta]}$ intersects $\Sigma$ eventually only at the regular locus $\Sigma_1$. We will prove that this implies that $(D\phi_{ t}L)_{t\,\in\,[\alpha, \beta]}$ is actually topologically transverse to $\Sigma$ and that the Maslov index is non-positive.

Let $t_0\in [\alpha, \beta]$ be such that $D\phi_{t_0}L\in \Sigma_1$. We introduce $x_0=p(L)$ and $x_t:=\phi_t(x_0)$.\linebreak Consider a continuous Lagrangian bundle $K$ defined in a neighbourhood of $x_{t_0}$ and transverse to the vertical bundle. By hypothesis, there exists $\varepsilon>0$ such that
\begin{itemize}
\item $\forall t\in (t_0, t_0+\varepsilon), Q_{K(x_t)}(D\phi_t\circ (D\phi_{t_0})^{-1}V(x_{t_0}), V(x_t))$ is positive definite;
\item $\forall t\in (t_0-\varepsilon, t_0), Q_{K(x_t)}(D\phi_t\circ (D\phi_{t_0})^{-1}V(x_{t_0}), V(x_t))$ is negative definite.
\end{itemize}
Then, for $t\in [t_0, t_0+\varepsilon)$, if $L_t=D\phi_tL$, we have
\begin{multline*}
\cq_{K(x_t)}(L_t, V(x_t))=
\cq_{K(x_t)}\left(D\left(\phi_t\circ \phi_{t_0}^{-1}\right)L_{t_0}, D\left(\phi_t\circ \phi_{t_0}^{-1}\right)V(x_{t_0})\right)\\+\cq_{K(x_t)}\left(D\left(\phi_t\circ \phi_{t_0}^{-1}\right)V(x_{t_0}), V(x_t)\right).
\end{multline*}
Then:
\begin{itemize}
\item because of the invariance under symplectic diffeomorphisms, the signature of $\cq_{K(x_t)}(D(\phi_t\circ \phi_{t_0}^{-1})L_{t_0}, D(\phi_t\circ \phi_{t_0}^{-1})V(x_{t_0}))$ is equal to the signature of $\cq_{D(\phi_{t_0}\circ \phi_{t}^{-1})^{-1}K(x_t)}(L_{t_0}, V(x_{t_0}))$ and if we chose $\varepsilon$ small enough, this signature is equal to the signature of $\cq_{K(x_{t_0})}(L_{t_0}, V(x_{t_0}))$;
\item the quadratic form
\[
\cq_{K(x_t)}\left(D\left(\phi_t\circ \phi_{t_0}^{-1}\right)V(x_{t_0}), V(x_t)\right)=-\cq_{K(x_t)}\left(V(x_t), D\left(\phi_t\circ \phi_{t_0}^{-1}\right)V(x_{t_0})\right)
\]
is negative definite because we assume that the isotopy $(\phi_t)$ strictly twists the vertical.
\end{itemize}
Since the index of the sum of a quadratic form $Q$ and a negative {definite} quadratic form is at least  the sum of the index and the nullity of $Q$, we deduce that for $t\in (t_0, t_0+\varepsilon)$ the index of $\cq_{K(x_t)}(L_t, V(x_t))$ is at least the  sum of $1$, which is the nullity of $\cq_{K(x_{t_0})}(L_{t_0},V(x_{t_0}))$, and the index of $\cq_{K(x_0)}(L_{t_0}, V(x_0))$. { Moreover, as the quadratic form $-\cq_{K(x_t)}(V(x_t), D(\phi_t\circ \phi_{t_0}^{-1})V(x_{t_0}))$ is close to $0$ for $\varepsilon$ small enough and because of the continuous dependence of the eigenvalues on the quadratic form, the index of $\cq_{K(x_t)}(L_t, V(x_t))$ is exactly the sum of $1$ and the index of $\cq_{K(x_0)}(L_{t_0}, V(x_0))$.} A similar argument gives that for $t\in (t_0-\varepsilon, t_0)$, the index of $\cq_{K(x_t)}(L_t, V(x_t))$ is {exactly} the index of $\cq_{K(x_0)}(L_{t_0}, V(x_0))$ This proves that $(L_t)_{t\,\in\,[\alpha, \beta]}$ intersect $\Sigma_1$ { topologically} transversally and in the negative sense and that the Maslov index is non-positive.\newline
Let now $(\phi_t)$ be an isotopy that twists the vertical (with no further assumptions). Consider the Lagrangian path $(D\phi_tL)_{t\,\in\,[\alpha,\beta]}$.

\begin{enonce}{Claim}
There exists an isotopy $(\tilde\phi_t)$ of conformally symplectic diffeomorphisms of $\M$ such that
\begin{itemize}
\item $(D\tilde\phi_tL)_{t\,\in\,[\alpha,\beta]}$ is a smooth perturbation of $(D\phi_tL)_{t\,\in\,[\alpha,\beta]}$, and in particular, $\MI((D\phi_tL)_{t\,\in\,[\alpha,\beta]})=\MI((D\tilde\phi_tL)_{t\,\in\,[\alpha,\beta]})$;
\item $(\tilde\phi_t)$ is an isotopy {that satisfies the conclusion of Proposition\ref{Ptwistintrinsec} (with definite quadratic forms) in a neighborhood of $(D\tilde\phi_tL)_{t\,\in\,[\alpha,\beta]}$.}
\end{itemize}
\end{enonce}

The claim immediately implies that the Maslov index of $(D\phi_tL)_{t\,\in\,[\alpha,\beta]}$ is non-positive, as desired.\newline
Let us now prove the claim. Because $(\phi_t)$ twists the vertical, we deduce from Proposition~\ref{Ptwistintrinsec} that there exists $\varepsilon>0$ such that
\begin{itemize}
\item $\forall t\in (t_0, t_0+\varepsilon), Q_{K(x_t)}(D\phi_t\circ (D\phi_{t_0})^{-1}V(x_{t_0}), V(x_t)) $ is a negative semi-definite quadratic form;
\item $\forall t\in (t_0-\varepsilon, t_0), Q_{K(x_t)}(D\phi_t\circ (D\phi_{t_0})^{-1}V(x_{t_0}), V(x_t)) $ is a positive semi-definite quadratic form.
\end{itemize}
If the vector field associated to $(\phi_t)$ is written in the chart as $X=(X_q, X_p)$, we deduce from equations~\eqref{Ehamlin} that
\[
\frac{d}{dt}\big(b_t(d_t)^{-1}
\big)=\partial_qX_qb_t(d_t)^{-1}+\partial_pX_q-b_t(d_t)^{-1}\partial_qX_pb_t(d_t)^{-1}-b_t(d_t)^{-1}\partial_pX_p\,.
\]
Since $b_t(d_t)^{-1}$ is the matrix of $Q_{K(x_t)}(D\phi_t\circ \big(D\phi_{t_0}\big)^{-1}V(x_{t_0}), V(x_t))$ and is zero for $t=t_0$, we deduce that
\[
\frac{d}{dt}\big(b_t(d_t)^{-1}
\big)_{|t=t_0}= \partial_pX_q\,.
\]
and then that $\partial_pX_q$ is a positive semi-definite quadratic form. We now add to $X$ a small Hamiltonian vector-field $Y$ that is associated to a Hamiltonian $H$ that is strictly convex in the fiber direction. This implies that $\partial_pY_q$ is positive definite and so is $\partial_p(X+Y)$. {According to Proposition~\ref{Ptermchpvect},} the isotopy associated with $X+Y$ is the desired isotopy $(\tilde\phi_t)$.
\end{proof}


\begin{prop}\label{PHconvextwist}
Assume that $a:I\rightarrow \R$ and $H:I\times T^*M\rightarrow \R$ are smooth functions and let us use the notation $H_t(x)=H(t, x)$. We assume that the Hessian of $H$ restricted to every vertical fiber is positive definite. We define the time-dependent vector field $X_t$ on $T^*M$ by
\[
i_{X_t}\omega=dH_t-a(t)\lambda.
\]
Then the isotopy defined by $X_t$ is conformally symplectic and strictly twists the vertical.
\end{prop}

\begin{proof}[Proof of Proposition~\ref{PHconvextwist}.] For $(t_0, x_0)\in I\times T^*M$, we choose a vertically foliated Darboux chart $F=(F_1, F_2):\cU \rightarrow \R^d\times \R^d$ such that $F(x_0)=0$ and $F(\cal U)=(-a, a)^{d}\times (-a, a)^d$.\newline
We now work in this chart and denote by $H$ the Hamiltonian in this chart, which has a positive definite Hessian in the $p$ direction. We chose $\varepsilon_0 >0$ such that for all $t\in (t_0-\varepsilon_0, t_0+\varepsilon_0)$, ${\cg}_t:=(\phi_t\circ\phi_{t_0}^{-1})(F^{-1}(\{ 0_{\R^d}\}\times (-\frac{a}{2}, \frac{a}{2})^d)) \subset\cU$ and $F({\cg_t})$\footnote{$F({\cg_t})$ is Lagrangian.} is the graph of a function $p\mapsto q=dg_t(p)$. We deduce from the Hamilton equations that there exists $\varepsilon\in (0, \varepsilon_0)$ such that uniformly for $y\in \ (-\frac{a}{2}, \frac{a}{2})^d$ and $t\in (t_0-\varepsilon, t_0+\varepsilon)\backslash \{t_0\}$ if we use the notation $\phi_t(0, y)=(q_t, p_t)$ then
\[
D^2g_t(p_t)= (t-t_0)\left(\frac{\partial^2 H}{\partial p^2}(t_0, 0, y)+O(t-t_0)\right).
\]
This results in the {(strict) }twist property.
\end{proof}

\subsection{Maslov index and symplectic reduction}\label{sym reduction}
On a cotangent bundle, the Maslov index is invariant by symplectic reduction. The result is due to C.~Viterbo~\cite{Vit1987}. For sake of completeness, we recall here Viterbo's proof.


Let us start by showing the invariance of the Maslov index by symplectic reduction on a symplectic vector space. Let $(V,\omega)$ be a symplectic vector space of dimension~$2d$. Denote by $\Lambda(V)$ the set of Lagrangian subspaces in $V$ and, for every subspace $U\subset V$, by $\Lambda_U(V)$ the set of Lagrangian subspaces $L$ such that $L\cap U=\{0\}$.

Fix $L_0\in\Lambda(V)$\footnote{$L_0$ is the Lagrangian subspace with respect to which we calculate the Maslov index in $(V,\omega)$.}. Let $W\subset V$ be a coisotropic (non-Lagrangian) vector subspace such that
\[
W^\perp\subset L_0\subset W\,,
\]
where $W^\perp$ denotes the symplectic orthogonal with respect to $\omega$. Consider the quotient map
\begin{equation*}
\begin{aligned}
\Pi_{W^\perp}:& \,W\to W/W^\perp\\
&\,v\mapsto [v]\,,
\end{aligned}
\end{equation*}
where $[v]=[u]$ if and only if $v-u\in W^\perp$. Observe that $\Pi_{W^\perp}$ is a surjective linear map. Then, the quotient space inherits a symplectic 2-form $\omega_W$ from $\omega$, and $(W/W^\perp,\omega_W)$ is still a symplectic vector space. In particular, for every Lagrangian subspace $L$ of $V$ the image $\Pi_{W^\perp}(L\cap W)$ is still a Lagrangian space in $W/W^\perp$. Denote by
\[
\P_{W^\perp}:\Lambda(V)\hookrightarrow \Lambda\left(W/W^\perp\right)
\]
\[
L\mapsto \Pi_{W^\perp}(L\cap W)\,.
\]



The following holds.

\begin{enonce}{Claim}
The map $\P_{W^\perp}$ restricted to $\Lambda_{W^\perp}(V)$ is a submersion.
\end{enonce}
\begin{proof}[Proof of the claim.]
Let us fix $L\in\Lambda_{W^\perp}(V)$ and let $L'\in\Lambda(V)$ be such that $W^\perp\subset L'\subset W$ and $L\cap L'=\{0\}$. The set $U=\{ \tilde L\in\Lambda(V) :\ \tilde L\cap L'=\{0\} \}$ is an open neighbourhood of $L$. If $\tilde L\in U$, then there exists a unique linear map $B=B_{\tilde L}: L\to L'$ such that
\[
\tilde L=\{ v+Bv; v\in L\}
\]
and $B$ satisfies the symmetry condition
\begin{equation}\label{Esym}
\forall \ell, \ell'\in L, \omega(\ell, B\ell')+\omega(B\ell, \ell')=0.
\end{equation}
Moreover, if $B:L\to L'$ satisfies the symmetry condition~\eqref{Esym}, then the set $\tilde L=\{ v+Bv; v\in L\}$ is a Lagrangian subspace of $V$ that is transverse to $L'$.\newline
We denote by $\cb$ the set of linear maps from $L$ to $L'$ that satisfy the symmetry condition~\eqref{Esym}. It is a finite dimensional vector space that is the image of the chart
\[
\tilde L\in U\mapsto B_{\tilde L}\in \cb.
\]
Similarly, if $\overline{L}=\P_{W^\bot}(L)$ and $\overline{L}'=\P_{W^\bot}(L')$, the set
\[
V=\left\{\tilde L\in \Lambda\left(W/W^\bot\right); \tilde L\cap \overline{L}'=\{ 0\}\right\}
\]
is an open neighbourhood of $\overline{L}$ in $\Lambda(W/W^\bot)$. The map that associate to every $\tilde L\in V$ the linear map $\overline B_{\tilde L}:\overline{L}\to \overline{L}'$ such that $\tilde L=\{\ell+\overline B_{\tilde L}\ell;\ \ell\in\overline L\}$ is a chart whose image is the finite dimensional vector space $\overline{\cb}$ of linear maps $\overline B:\overline{L}\to \overline{L}'$ such that
\[
\forall \ell, \ell'\in \overline{L},\omega_W(\ell, \overline B\ell')+\omega_W(\overline B\ell, \ell')=0.
\]
In these charts, the map $\P_{W^\bot}$ is read $\Phi: \cb\to\overline{\cb} $ where
\[
\Phi(B)=\Pi_{W^\perp}\circ B|_{L\cap W}\circ \left(\Pi_{W^\perp}\,\middle|_{L\,\cap\,W}\right)^{-1}.
\]
Hence $\Phi$ is a linear map. This is then a submersion onto its image that is a linear subspace of $\overline{\cb}$. If we prove that $\Phi(\cb)=\overline{\cb}$, we will deduce that $\P_{W^\bot}$ is a submersion. Thus, let $\overline{B}_0\in \overline{\cb}$ and let $\overline{L}_0$ be the graph of $\overline{B}$. We choose a linear subspace $L'_1$ of $L'$ that is transverse to $W^\bot$ and define $B_1:L\cap W\to L'_1$ as
\[
\forall v\in L\cap W, B_1(v)=\Pi_{W^\perp}|_{L'_1}^{-1}\circ \overline{B}_0\circ \Pi_{W^\perp}(v).
\]
When $v, w\in L\cap W$, we have
\begin{align*}
\omega(v, B_1w)+\omega(B_1v, w)&=
\omega\left(v,\Pi_{W^\perp}\middle|_{L'_1}^{-1}\circ \overline{B}_0[w]\right) +\omega\left(\Pi_{W^\perp}\middle|_{L'_1}^{-1}\circ \overline{B}_0[v],w\right)\\
&=\omega_W\left([v], \overline{B}_0[w]\right) +\omega_W\left(\overline{B}_0[v],[w]\right)=0\,.\\
\end{align*}
Hence $B_1:L\cap W\to L'_1$ satisfies the symmetry condition and then its graph $L_2$ is an isotropic subspace of $(L\cap W)+L'\subset W$ such that $L_2\cap L'=\{ 0\}$. We can choose a Lagrangian subspace $\tilde L$ of $V$ that contains $L_2$ and is transverse to $L'$: $\tilde L$ is then the graph of a map $B_2:L\to L'$ that satisfies the symmetry condition and contains $L_2$. Then the graph of $\Phi(B_2)$ is a Lagrangian subspace of $W/W^\bot$ that contains $\overline{L}_0$, hence is equal to $\overline{L}_0$, from which we deduce that $\Phi(B_2)=\overline{B}_0$, demonstrating the surjectivity of $\Phi$.
\end{proof}


Denote by $i:\Lambda_{W^\perp}(V)\hookrightarrow\Lambda(V)$ the standard inclusion. This is a submersion.

\begin{lemm}\label{sym red for vec space}
Let $t\in [0,1]\mapsto \gamma(t)\in\Lambda_{W^\perp}(V)$ be an arc such that $\gamma(0)\cap L_0=\gamma(1)\cap L_0=\{0\}$. Then
\[
\MI(i\circ\gamma)=\MI\left(\P_{W^\perp}\circ\gamma\right)\,,
\]
where
\begin{itemize}
\item the Maslov index $ \MI(i\circ\gamma)$ is calculated with respect to $L_0$ in $\Lambda(W)$;
\item the Maslov index $ \MI(\P_{W^\perp}\circ\gamma)$ is calculated with respect to $\P_{W^\perp}(L_0)$ in $\Lambda(W/W^\perp)$.
\end{itemize}
\end{lemm}


\begin{proof}
By slightly perturbing the path, we can assume that $\gamma$ is in D-general position with respect to $\Sigma:=\{L\in\Lambda(V) ;\ L\cap L_0\neq \{0\}\}$. The subspace $L_0':= \P_{W^\perp}(L_0)$ of $W/W^\perp$ is Lagrangian and we denote $\overline{\Sigma}=\{L'\in\Lambda(W/W^\perp) ;\ L'\cap L_0'\neq\{0\}\}$. Observe that $\P_{W^\perp}^{-1}(\overline{\Sigma})\subset\Sigma$ because $W^\bot\subset L_0$.


Since the maps $i$ and $\P_{W^\perp}$ are submersions, we have
\begin{itemize}
\item the path $i\circ\gamma$ is in D-general position with respect to $\Sigma$;
\item the path $\P_{W^\perp}\circ\gamma$ is in D-general position with respect to $\overline{\Sigma}$.
\end{itemize}


Moreover, the choice of a coorientation of $\Sigma$ determines a coorientation both on $i^{-1}(\Sigma_1)$ and on $\P_{W^\perp}\circ i^{-1}(\Sigma_1)$. Following~\cite{Vit1987}, we claim that

\begin{enonce}{Claim}\label{claim23}
\[
i^{-1}({\Sigma)=
\P_{W^\perp}^{-1}(\overline{\Sigma}) \cap \Lambda_{W^\bot}(V)}\,.
\]
\end{enonce}

\begin{proof}[Proof of the claim.]We first observe that on one side
\[
i^{-1}({\Sigma)}=
\left\{ L\in\Lambda(V) :\ L\cap W^\perp=\{0\}\text{ and } L\cap L_0\neq\{0\} \right\}\,.
\]
On the other side we have
\begin{multline*}
\P_{W^\perp}^{-1}\left({\overline{\Sigma}}\right){\cap \Lambda_{W^\bot}(V)}=
\\
\left\{ L\in{ \Lambda_{W^\bot}(V)}\, ;\, \left((L\cap W+W^\perp)/W^\perp\right)\cap \left((L_0\cap W+W^\perp)/W^\perp\right)\neq \{0\} \right\}\, ;
\end{multline*}
since $W^\perp\subset L_0\subset W$, we have $L_0=L_0\cap W+W^\perp$ and any $L\in \P_{W^\perp}^{-1}({\overline{\Sigma}}){\cap \Lambda_{W^\bot}(V)}$ is so that\\
$ ((L\cap W+W^\perp)/W^\perp)\cap ({L_0 }/W^\perp)=$\hglue 5 truecm
\[
\left(L\cap L_0\cap W+W^\perp\right)/W^\perp=\left(L\cap L_0+W^\perp\right)/W^\perp\neq \{0\}\,.
\]
Since $L\cap W^\perp=\{0\}$
\[
\left(L\cap L_0+W^\perp\right)/W^\perp\neq \{0\}\quad\Leftrightarrow\quad L\cap L_0\neq \{0\}\,.
\]
We so conclude that
\begin{align*}
\P_{W^\perp}^{-1}({\overline{\Sigma}}){\cap \Lambda_{W^\bot}(V)}&=
\left\{ L\in \Lambda(V) :\ L\cap W^\perp=\{0\}\text{ and } L\cap L_0\neq \{0\}
\right\}\\
&=i^{-1}\left(\left\{ L\in\Lambda(V) :\ L\cap L_0\neq\{0\}\right\}\right)\,.
\end{align*}
\end{proof}


Since $i$ is a submersion, the number of crossings of $\gamma$ with $ i^{-1}({ \Sigma})$ is equal to the number of crossings of $i\circ\gamma$ with ${ \Sigma}$. {Since $\P_{W^\perp}$ is also a submersion and $\P_{W^\perp}(\Sigma)=\bar\Sigma$}, we conclude that the number of crossings of $\gamma$ with $ i^{-1}({ \Sigma})$ is equal to the number of crossings of $\P_{W^\perp}\circ\gamma$ with ${ \overline{\Sigma}}$.

Since the coorientation on $i^{-1}(\Sigma)$ and on $\P_{W^\perp}\circ i^{-1}(\Sigma)$ is determined by the coorientation of $\Sigma$, we conclude that actually the number of positive (resp. negative) crossings of $i\circ\gamma$ corresponds to the number of positive (resp. negative) crossings of $\P_{W^\perp}\circ\gamma$. By the definition of Maslov index, we obtain the sought result.
\end{proof}

We want now to prove the invariance of the Maslov index by symplectic reduction on the cotangent bundle $\M$, endowed with the symplectic form $\omega$. Let $\cv$ be the Lagrangian foliation whose fibers of the associated tangent Lagrangian bundle are the vertical Lagrangian subspaces. Let $\cal W\subset \M$ be a coisotropic submanifold {and let $i_\cw:\cal W\to\M$ be the canonical injection}; the characteristic foliation of $\cal W$,  denoted by $\cal W^\bot$, admits for tangent bundle ${T_x(\cal W^\bot)=}\ker(i^*_\cw\omega)(x)=(T_x\cal W)^\perp$.


Assume that, for every $x\in { \cal W}$ it holds
\[
(T_x\cal W)^\perp\subset T_x\cv\subset T_x\cal W\,.
\]
We assume that the symplectic reduction of $\cal W$ is a true symplectic manifold that we denote by $\cal R: \cal W\to \cal W/\cal W^\bot$. When $x\in\cal W$ and $L\in \Lambda_{T_x\cal W^\bot}(T_x\M)$, we denote
\[
\P(L)=D\cal R(x)L=\left(L+T_x\cal W^\bot\right)/T_x\cal W^\bot\in \Lambda\left(\cal W/\cal W^\bot\right).
\]
Then $\P$ is a submersion from $\Lambda_{\cal W^\bot}(\M)_{|\cal W}$ to $\Lambda(\cw/\cal W^\bot)$.



We denote $\Sigma(\M):= \{ L\in \Lambda(\M) ;\ L\cap T_{p(L)}\cv\neq \{0\} \}$ and $\Sigma(\cal W/\cal W^\bot)=\{ L\in \Lambda(\cal W/\cal W^\bot); L\cap T_{p(L)}\P(\cal V)\neq \{ 0\}\}$ .


\break
\begin{lemm}\label{lemma almost Viterbo}
Let $\Gamma:[a,b]\to \Lambda(\M)$ be a smooth arc such that
\begin{itemize}
\item $\Gamma(a),\Gamma(b)\notin \Sigma(\M)$;
\item {$\Gamma$ is in $D$-general position with respect to the fibered singular cycle $\Sigma(\M)$;}
\item at every point the path has trivial intersection with the tangent bundle of the characteristic foliation of $\cal W$, i.e.
\[
\Gamma(t)\cap \left(T_{p(\Gamma(t))}\cal W\right)^\perp=\{0\}\quad\forall t\in[a,b]\,.
\]
\end{itemize}
Then
\[
\MI(\Gamma)= \MI(\P\circ\Gamma)\,,
\]
where
\begin{itemize}
\item the Maslov index $ \MI(\Gamma)$ is calculated with respect to $T\cal V$ in $\Lambda(\M)$;
\item the Maslov index $ \MI(\P\circ\Gamma)$ is calculated with respect to $\P(\cal V)$ in $
\Lambda(\cal W/\cal W^\bot)$.
\end{itemize}
\end{lemm}

\begin{proof}
Since $\P$ is a submersion and $\P^{-1}(\Sigma(\cal W/\cal W^\bot))\cap \Lambda_{T\cal W^\bot}(\M)=\Sigma(\M)_{|\cal W}$, the path $\P\circ\Gamma$ is also in $D$-general position with respect to $\Sigma(\cal W/\cal W^\bot)$. In order to conclude, it is then sufficient to calculate the Maslov index of a sub-path of $\Gamma$, around a (isolated) crossing $t$. Let $U\subset\M$ be a neighborhood of $p\circ\Gamma(t)$ and let
\[
\Gamma\vert_{[t-\epsilon,t+\epsilon]}:[t-\epsilon,t+\epsilon]\to \Lambda(U)\,,
\]
be a Lagrangian path with only an isolated, transverse crossing at $t$. Let us trivialise $\Lambda(U)$ as $U\times \Lambda(T_{p\circ\Gamma(t)}\M)$. Similarly, trivialise the image $\P(\Lambda(U))$ as $\P(U)\times \Lambda(T_{p\circ\Gamma(t)}\cw/(T_{p\circ\Gamma(t)}\cw)^\perp)$. Up to restrict the neighborhood $U$, the Maslov index of the path $\Gamma\vert_{[t-\epsilon,t+\epsilon]}$ with respect to $\cv$ corresponds to the Maslov index of $\Gamma\vert_{[t-\epsilon,t+\epsilon]}$, seen as a Lagrangian path in the symplectic vector space $T_{p\circ\Gamma(t)}\M$ thanks to the trivialization, with respect to $T_{p\circ\Gamma(t)}\cv$. Similarly, the Maslov index of the path $\P(\Gamma\vert_{[t-\epsilon,t+\epsilon]})$ with respect to $\P(\cv)$ is actually the Maslov index of the Lagrangian path $\P(\Gamma\vert_{[t-\epsilon,t+\epsilon]})$, seen in $T_{p\circ\Gamma(t)}\cw/(T_{p\circ\Gamma(t)}\cw)^\perp$ through the trivialization, with respect to $\P(T_{p\circ\Gamma(t)}\cv)$. By applying Lemma~\ref{sym red for vec space}, we can then conclude.
\end{proof}






\section{Maslov index along a Lagrangian submanifold that admits a generating function}\label{section gfqi}
Let $\cl\subset T^*M$ be a Lagrangian submanifold. The goal of this Section is to prove that every arc $\Gamma:[a,b]\to T\cl$ whose endpoints project on $T^*M$ on the so-called graph selector of $\cl$ has zero Maslov index.
\subsection{The relation between the Maslov index and the Morse index} Let us recall the definition of generating function for a Lagrangian submanifold $\cl$ of $T^*M$.

\begin{defi}
A $C^r$ function with $r\geq 2$ function $S:M\times \R^k\rightarrow \R$ {\sl generates} a Lagrangian submanifold $\cl$ of $T^*M$ if
\begin{itemize}
\item using the notation
\[
\cC_S=\left\{ (q, \xi)\in M\times \R^k :\ \frac{\partial S}{\partial \xi}(q, \xi)=0\right\},
\]
at every point of $\cC_S$, the map $\frac{\partial S}{\partial \xi}$ is a submersion; in this case, $\cC_S$ is a $d$-dimensional submanifold of $M\times \R^k$;
\item the map $j_S: \cC_S\hookrightarrow T^*M$ defined by $j_S(q, \xi)=\frac{\partial S}{\partial q}(q, \xi)$ is an embedding such that $j_S(\cC_S)=\cl$.
\end{itemize}
The \emph{generating function} $S$ is \emph{quadratic at infinity} (GFQI) if there exists a compact subset $K\subset M\times \R^k$ and a non-degenerate quadratic form $Q:\R^k\rightarrow \R$ such that
\[
\forall (q, \xi)\notin K, S(q, \xi)=Q(\xi).
\]
The generating function quadratic at infinity $S$ is of index $m$ if the non-degenerate quadratic form $Q$ has index $m$.
\end{defi}

A result due to Sikorav~\cite{Bru1991,Sik1987} asserts that every H-isotopic\footnote{This means Hamiltonianly isotopic} to the zero section submanifold of $T^*M$ admits a GFQI.

\begin{nota}
If we denote as before the Liouville form on $T^*M$ by $\lambda$ and the Liouville form on $T^*(\R^k)$ by $\lambda_1$, the product manifold $\cN=T^*M\times T^*(\R^k)$ is endowed with the symplectic form
$\Omega=-\bbP_1^*d\lambda-\bbP_2^*d\lambda_1$ where $\bbP_i$ is the projection on the i$^{\rm th}$ factor.
\end{nota}

\begin{theo}\label{Tmain}
Let $\cl\subset T^*M$ be a Lagrangian submanifold that admits a generating function $S(q, \xi):M\times \R^k\rightarrow \R$. Let $(q_i, \xi_i)\in M\times \R^k$, $i= 1, 2$ be such that
\begin{itemize}
\item $\frac{\partial S}{\partial \xi}(q_i, \xi_i)=0$, i.e. $(q_i,\xi_i)\in \cC_S$;
\item if we use the notation $p_i= \frac{\partial S}{\partial q}(q_i, \xi_i)$, the submanifold $\cl$ is transverse to the vertical fiber $T^*_{q_i}M$ at $p_i$ in $T^*M$.
\end{itemize}
Then, $\ker \frac{\partial^2 S}{\partial \xi^2}(q_i, \xi_i)=\{ 0\}$ and for every arc $\gamma_0$ joining $\gamma_0(0)=p_1$ to $\gamma_0 (1)=p_2$ in $\cl$, the Maslov index of $t\in [0, 1]\mapsto T_{\gamma_0 (t)}\cl$ with respect to the vertical is equal to the difference of the Morse indices $\text{index}(\frac{\partial^2 S}{\partial \xi^2}(q_2, \xi_2))-\text{index}(\frac{\partial^2 S}{\partial \xi^2}(q_1, \xi_1))$.
\end{theo}

\begin{proof}[Proof of Theorem~\ref{Tmain}]

\begin{lemm}\label{Lcaractvert}
Let $p=\frac{\partial S}{\partial q}(q, \xi)\in \cl$. Then $\cl$ is transverse to $T^*_qM$ at $p$ if and only if$ker (\frac{\partial^2 S}{\partial \xi^2}(q, \xi))=\{ 0\}$.
\end{lemm}
\begin{proof}[Proof of Lemma~\ref{Lcaractvert}]
Let us fix $p\in \cl$ and let $\delta p\in T_p(T^*M)$. We use the notation $q=\pi(p)\in M$ and $\delta q= d\pi(p)\delta p\in T_qM$.\\
Then $\delta p$ belongs to $T_{p}\cl$ if and only if there exists $\delta\xi\in \R^k$ such that
\begin{itemize}
\item ${D(\frac{\partial S}{\partial \xi})(\delta q, \delta \xi)=}\frac{\partial^2 S}{\partial q\partial \xi}(q, \xi)\delta q+\frac{\partial^2 S}{\partial \xi^2}(q, \xi)\delta \xi=0$;
\item ${ Dj_S(\delta q, \delta \xi)=}\delta p=\frac{\partial^2 S}{\partial q^2}(q, \xi)\delta q+\frac{\partial^2 S}{\partial \xi\partial q}(q, \xi)\delta \xi$.
\end{itemize}
Observe that $\pi(\frac{\partial S}{\partial q}(q, \xi))=q$ and then
\begin{equation}\label{Egenevert}
D\pi\left(\frac{\partial^2S}{\partial q^2}(q, \xi)\delta q\right)=\delta q\quad\text{and}\quad D\pi\left(\frac{\partial^2S}{\partial \xi\partial q}(q, \xi)\delta \xi\right)=0.
\end{equation}
We deduce
\begin{gather*}
\delta\xi\in\ker \left(\frac{\partial^2S}{\partial\xi^2}\right)\backslash\{ 0\}\Longleftrightarrow(0, \delta \xi)\in \ker\left(D\left(\frac{\partial S}{\partial \xi}\right)\right)\backslash\{ 0\}
\\
\Longleftrightarrow Dj_S(0, \delta \xi)\in T\cl\backslash\{ 0\}\Longleftrightarrow\frac{\partial^2S}{\partial \xi\partial q}\delta\xi\in T\cl\backslash\{0\}.
\end{gather*}
Using~\eqref{Egenevert}, we conclude that $\ker (\frac{\partial^2 S}{\partial \xi^2}(q, \xi))\neq\{ 0\}$ if and only if $\cl$ is not transverse to $T^*_qM$ at $\frac{\partial S}{\partial q}(q, \xi)$.
\end{proof}

In ${\cal N}=T^*M\times T^*(\R^k)$, endowed with the symplectic form $\Omega${$=-\bbP_1^*d\lambda-\bbP_2^*d\lambda_1$}, we consider the coisotropic foliation into submanifolds
\[
\cw_\chi=T^*M\times \R^k\times \{\chi\}
\]
for $\chi\in \R^k$. The characteristic leaves of $\cw_\chi$ are the submanifolds $\cw^\bot_{(p,\chi)}=\{ p\}\times \R^k\times \{\chi\}$ with $ p\in T^*M$.



We will use also the Lagrangian foliation $\cF$ of $\cN$ with leaves $\cF_{q,\chi}=T^*_qM\times \R^k\times \{ \chi\}$. Then we have $\cw^\bot_{(p,\chi)} \subset \cF_{(\pi(p),\chi)}\subset \cw_\chi$. We denote by $F_{(p,\xi, \chi)}$ the tangent space to the leaf $\cF_{(\pi(p),\chi)}$ at the point $(p,\xi, \chi)$.\\
The graph $ \cg= \Graph(dS)\subset \cN$ of $dS$ is a Lagrangian submanifold of $\cN$ that is transverse to $\cw_0$ and such that $\cg\cap \cw_0$ is diffeomorphic to $\cl$ by the map
\[
\cR:(p, \xi, 0)\in \cw_0\mapsto p.
\]
Observe that $\cR$ is the symplectic reduction of $\cw_0$. We denote by $R$ the restriction of $\cR$ to $\cg\cap \cw_0$.

We use for $\gamma_0$, $q_i, p_i, \xi_i$ the same notations as in Theorem~\ref{Tmain}. Then $\Gamma_0=R^{-1}\circ \gamma_0$ is an arc on $\cg\cap \cw_0$ such that $\Gamma_0(0)=(p_1, \xi_1, 0)$ and $\Gamma_0 (1)=(p_2, \xi_2, 0)$. We have

\begin{lemm}\label{Lindices}
Let $\Gamma(t)=(p(t), \xi(t), \chi(t))\in \cg$ be an arc in $\cg$ such that at $\Gamma(0)$ and $\Gamma(1)$, the quadratic form $\frac{\partial^2 S}{\partial \xi^2}$ is non-degenerate. The Maslov index of the arc of Lagrangian subspaces $t\in[0, 1]\mapsto T_{\Gamma(t)}\cg$ with respect to the fibered singular cycle associated to $\cF$
is
\[
\text{index}\left(\frac{\partial^2 S}{\partial \xi^2}(q(1), \xi(1))\right)-\text{index}\left(\frac{\partial^2 S}{\partial \xi^2}(q(0), \xi(0))\right).
\]
\end{lemm}


\begin{proof}[Proof of Lemma~\ref{Lindices}]
Up to a small perturbation, we can assume that $S$ is smooth. The proof is divided into two steps. First, we will perturb the Lagrangian submanifold $\cl$ (i.e., its generating function) and $\Gamma$ on it so that $\Gamma$ is in $D$-general position with respect to the fibered singular cycle associated to $\cF$. Then, we will prove the lemma. 
%\noindent\emph{First step.}

\begin{step}\label{step1}
As $T_{\Gamma(0)}\cg$ and $T_{\Gamma(1)}\cg$ are transverse to $F_{\Gamma(0)},F_{\Gamma(1)}$ respectively, there exists $\varepsilon>0$ such that for all $ t\in [0, \varepsilon]\cup[1-\varepsilon, 1], T_{\Gamma(t)}\cg$ is transverse to $F_{\Gamma(t)}$. We use the notation $t\mapsto \zeta(t)$ $ := (\pi\circ p(t), \xi(t))\in M\times \R^k$. We now choose a neighbourhood $\cal U$ of {$\zeta([\varepsilon, 1-\varepsilon])$} in $M\times \R^k$ and a diffeomorphism $\psi: \cal U\rightarrow \R^d\times \R^k$ such that
\[
\forall t\in [\epsilon,1-\epsilon],
{\psi(\zeta (t))}=(t, 0,\, \dots,\, 0).
\]
Let $s:\psi(\cU)\rightarrow \R$ defined by $s(y)=S\circ \psi^{-1}(y)$. Then in the neighborhood of $(\varepsilon, 0\dots,\,0)$ and $(1-\varepsilon, 0,\,\dots,\,0)$, we know that $\Graph(D^2s)$ and $D\psi(F)$ are transverse. We can slightly perturb the path of matrices $t\in [\varepsilon, 1-\varepsilon]\mapsto D^2s(t, 0,\,\dots,\,0)$ in a path $t\mapsto A(t)$ of symmetric matrices such that
\begin{itemize}
\item $A(t)=D^2s(t, 0,\,\dots,\,0)$ in a neighbourhood of $\varepsilon$ and $1-\varepsilon$;
\item the path $t\mapsto \Graph(A(t))$ is in $D$-general position
\end{itemize}
We now define for $x_1$ in a neighborhood of $[\varepsilon, 1-\varepsilon]$
\begin{itemize}
\item $\delta(x_1)=\int_\varepsilon^{x_1}(A(\sigma)-D^2s(\sigma, 0,\,\dots,\,0))(1, 0,\, \dots,\, 0)d\sigma$;
\item $v(x_1)=\int_\varepsilon ^{x_1}\delta(\sigma)(1,0,\,\dots,\,0)d\sigma$;
\item in a neighbourhood of $[\varepsilon, 1-\varepsilon]\times \{ 0_{\R^{d-1}\times \R^k}\}$,
\begin{multline*}
u\left(x_1,\,\dots,\,x_{d+k}\right)=s\left(x_1,\,\dots,\,x_{d+k}\right)+v(x_1)+\delta(x_1)\left(0,x_2, \dots, x_{d+k}\right)
\\
+\frac{1}{2}(A(x_1)-D^2s(x_1,0,\,\dots,\,0))\big((0, x_2, \,\dots,\,x_{d+k}), (0, x_2, \,\dots,\, x_{d+k})\big).
\end{multline*}
\end{itemize}
Then $u$ is $C^2$ close to $s$ and we have
\[
\forall x_1\in [\varepsilon, 1-\varepsilon], D^2u(x_1, 0,\,\dots,\,0)=A(x_1).
\]
We then use a bump function $\eta$ with support in a neighbourhood of $[\varepsilon, 1-\varepsilon]\times \{ 0_{\R^{d-1}\times \R^k}\}$ and that is equal to 1 in a smaller neighbourhood of $[\varepsilon, 1-\varepsilon]\times \{ 0_{\R^{d-1}\times \R^k}\}$. We define
\begin{multline*}
\tilde s(x_1,\,\dots,\,x_{d+k})=
\\
(1-\eta(x_1, \,\dots,\,x_{d+k}))s(x_1, \,\dots,\, x_{d+k})+\eta(x_1, \,\dots,\, x_{d+k}) u(x_1, \,\dots,\, x_{d+k}).
\end{multline*}
As $\tilde s$ is equal to $u$ in $[\varepsilon, 1-\varepsilon]\times \{ 0_{\R^{d-1}\times \R^k}\}$, $D^2\tilde s$ is in {$D$-general position} with respect to $D\psi(F)$ along the lift of this arc in $\Graph\,Du$. In addition, as $\tilde s$ is $C^2$-close to $s$, $D^2\tilde s$ is transverse to $D\psi(F)$ along the lift of $([0, \varepsilon]\cup[1-\varepsilon, 1])\times \{ 0_{\R^{d-1}\times \R^k}\}$ in $\Graph\,D\tilde s$.

Finally, define the function $\tilde S$ to be equal to $S$ outside $\psi^{-1}(\cU)$ and to $\tilde s\circ \psi$ in $\psi^{-1}(\cU)$. Thus, $\tilde S$ is $C^2$ close to $S$, {$D\tilde S\circ \zeta$ is $C^1$ close to $DS\circ\zeta=\Gamma$ and $t\mapsto D\tilde S\circ\zeta(t)$ is in {$D$-general position} with respect to the fibered singular cycle associated to $\cF$}. As the new generating function $\tilde S$ is $C^2$ close to $S$, the number $\text{index}(\frac{\partial^2 S}{\partial \xi^2}(q(1), \xi(1)))-\text{index}(\frac{\partial^2 S}{\partial \xi^2}(q(0), \xi(0)))$ does not change. 


This will allow us to assume that the path $t\mapsto T_{\Gamma(t)}\cg$ is in {$D$-general position} with respect to {the fibered singular cycle associated to $\cF$}.
\end{step}

%\noindent \emph{Second step.}
\begin{step}\label{step2}
We then look at what happens at a crossing $\Gamma(\bar t)$. We choose a chart close to $\frac{\partial S}{\partial q}(q(\bar t), \xi(\bar t))$ and we assume that we work in coordinates~: $q\in U\subset \R^n$ and $(p_1,\, \dots,\,p_n)$ are the dual coordinates defined by $(q, \sum p_idq_i)\in T^*U$.\newline
In these coordinates for $(q, p)=(q, \frac{\partial S}{\partial q}(q, \xi))\in U\times \R^n$, we use the linear and symplectic change of coordinates

\[
(\delta q, \delta p)\mapsto \left(\delta Q=\delta p+\left(\bfone_n-\frac{\partial^2 S}{\partial q^2}(q, \xi)\right)\delta q, \delta P=-\delta q\right).
\]
In the extended space $T\cN$ with linear coordinates $(\delta Q, \delta P, \delta \xi, \delta \chi)$, the equation of $F_{\Gamma(t)}$ is $(\delta P, \delta\chi)=(0, 0)$, i.e., $F_{\Gamma(t)}$ is the graph of the zero function. The equation of $T\cg$ is
\begin{equation}\label{ETG}
\begin{cases}\delta P=-\delta Q+\frac{\partial^2S}{\partial \xi\partial q}\delta \xi\\
\delta \chi=\frac{\partial^2 S}{\partial q\partial \xi} \delta Q+\left(\frac{\partial ^2S}{\partial \xi^2} -\frac{\partial^2 S}{\partial q\partial \xi} \frac{\partial^2S}{\partial\xi\partial q} \right)\delta \xi
\end{cases}
\end{equation}
and this is also a graph. We then compute the change of Maslov index with respect to $F_{\Gamma(t)}$ with the help of the height of $T\cg$ above the vertical $L$ with respect to $F$, {i.e., $\cal Q_F(L,T\cg)$,} where $L$ has equation $(\delta Q, \delta\xi)=(0, 0)$. {Observe that, for $t$ close to $\bar t$}, $F$ and $L$ are transverse and the projection $\bbP_F:T\cN \rightarrow T\cN/F$ restricted to $L$ is an isomorphism. Therefore, we can take $(\delta P, \delta \chi)$ as coordinates in $T\cN/F$. Additionally, for $t\not=\bar t$ close to $\bar t$, $F$ and $T\cg$ are transverse, because crossings of a path in general position are isolated. Moreover, note that, for $t$ close to $\bar t$, $L$ and $T\cg$ are transverse. If we introduce the matrix
\begin{equation}
M(t)=
\begin{pmatrix} -\bfone_n & \frac{\partial^2 S}{\partial\xi\partial q}\\
\frac{\partial^2 S}{\partial q\partial \xi} & \left(\frac{\partial ^2S}{\partial \xi^2}-\frac{\partial^2 S}{\partial q\partial \xi}\frac{\partial^2S}{\partial\xi\partial q}\right)
\end{pmatrix}(q(t), \xi(t)),
\end{equation}
as the equation of $T\cg$ is (we write in coordinates)
\[
\begin{pmatrix}
\delta P \\ \delta \chi
\end{pmatrix} = M(t)\,
\begin{pmatrix}
\delta Q \\ \delta \xi
\end{pmatrix}\,,
\]
the matrix $M(t)$ is invertible for $t\not=\bar t$ and we have
\[
\cq_F(L, T\cg)(\delta P, \delta\chi)=\Omega\Big((0,0,\delta P, \delta \chi), \left((\delta P,\delta\chi).\left(M(t)^{-1}\right)^T, \delta P, \delta\chi\right)\Big).
\]

Hence the matrix of $\cq_F(L, T\cg)(\delta P, \delta\chi)$ in coordinates $(\delta P, \delta\chi)$ is $M(t)^{-1}$. The change of signature of $M(t)^{-1}$ at $\bar t$ is exactly the same as the change of signature of $M(t)$.

Let us introduce the matrix
\[
\P(t)=
\begin{pmatrix} \bfone_n& \frac{\partial^2 S}{\partial\xi\partial q}(q(t), \xi(t))\\ {\bf 0}& \bfone_k
\end{pmatrix}.
\]
Then we have
\[
\begin{matrix}\P(t)^T M(t)\P(t)&=
\begin{pmatrix} \bfone_n& {\bf 0}\\ \frac{\partial^2 S}{\partial q\partial \xi}(q(t), \xi(t))& \bfone_k
\end{pmatrix}M(t)
\begin{pmatrix} \bfone_n& \frac{\partial^2 S}{\partial\xi\partial q}(q(t), \xi(t))\\ {\bf 0}& \bfone_k
\end{pmatrix}\\
&=
\begin{pmatrix} -\bfone_n & 0\\ 0&\frac{\partial ^2S}{\partial \xi^2}(q(t), \xi(t))
\end{pmatrix}\,.\hfill\
\end{matrix}
\]
Hence the change of signature of $M(t)$ at $t=\bar t$ along the path $\Gamma$ is equal to the change of signature of $\frac{\partial^2 S}{\partial\xi^2}$. This is exactly the Maslov index of the arc of Lagrangian subspaces $t\in[\bar t-\varepsilon, \bar t+\varepsilon]\mapsto T_{\Gamma(t)}\cg$ with respect to $F_{\Gamma(t)}$.
\end{step}
\end{proof}

We deduce that the Maslov index of $T\cg$ along the arc $\Gamma_0$ with respect to $F$ is $\text{index}(\frac{\partial^2 S}{\partial \xi^2}(q_2, \xi_2))-\text{index}(\frac{\partial^2 S}{\partial \xi^2}(q_1, \xi_1))$.\\
We have noticed that $\cw^\bot_{(p,\chi)} \subset \cF_{(\pi(p),\chi)}\subset \cw_\chi$. Also, because $\cg$ is Lagrangian and transverse to $\cw_0$, at every point of intersection, the intersection of the tangent subspaces to $\cg$ and $\cw^\bot_{(p,0)}$ is $\{ 0\}$. {The path $t\mapsto \Gamma(t)$ can be put in $D$-general position with respect to $\cF$, as done in the first step of the proof of Lemma~\ref{Lindices}. Thus, we can apply the results concerning the Maslov index that are given in~\cite[Section~2]{Vit1987}, see here Lemma~\ref{lemma almost Viterbo} and Subsection~\ref{sym reduction}}.
\\
As the curve $\Gamma_0$ is contained in $\cg\cap \cw_0$, the Maslov index of $T\cg$ along $\Gamma_0$ with respect $F$ is equal to the Maslov index of $(T(\cg\cap \cw_0))/T\cw^\bot_0$ with respect to $F/T\cw^\bot_0$. We have $(T(\cg\cap \cw_0))/T\cw^\bot_0={ T}R(\cg)={ T}\cl$ and $F/T\cw^\bot_0(\Gamma(t))=T_{\gamma_0(t))}(T^*_{\pi\circ \gamma_0(t)}M)$ is the vertical {$V_{\gamma_0(t)}$}. This proves the Theorem~\ref{Tmain}.
\end{proof}

\begin{thebibliography}{99}
\bibitem[MS16]{McDuffSalamon2017} D. McDuff and D. Salamon, \emph{Introduction to symplectic topology}, third ed., Oxford Graduate Texts in Mathematics, vol. 27, Oxford University Press, 2016.
\bibitem[RS93]{RobbinSalamon1993} J. Robbin and D. Salamon, ``The Maslov index for paths'', \emph{Topology} \textbf{32} (1993), no. 4, 827--844.
\bibitem[Dui76]{Dui76} J. J. Duistermaat, ``On the Morse index in variational calculus'', \emph{Adv. Math.} \textbf{21} (1976), no. 2, 173--195.
\bibitem[Arn08]{Arn08} M.-C. Arnaud, ``Fibr\'es de Green et r\'egularit\'e des graphes $C^0$-lagrangiens invariants par un flot de Tonelli'', \emph{Ann. Henri Poincar\'e} \textbf{9} (2008), no. 5, 881--926.
\bibitem[Wei71]{Weinstein1971} A. Weinstein, ``Symplectic manifolds and their Lagrangian submanifolds'', \emph{Adv. Math.} \textbf{6} (1971), 329--346.
\bibitem[MBA72]{Arn72} V. P. Maslov, V. C. Bouslaev, and V. I. Arnol\'d, \emph{Th\'eorie des perturbations et m\'ethodes asymptotiques}, \'Etudes math\'ematiques, Paris: Dunod; Gauthier-Villars, 1972.
\bibitem[CLM94]{Capleem} S. E. Cappell, R. Lee, and E. Y. Miller, ``On the Maslov index'', \emph{Commun. Pure Appl. Math.} \textbf{47} (1994), no. 2, 121--186.
\bibitem[Vit87]{Vit1987} C. Viterbo, ``Intersection de sous-vari\'et\'es lagrangiennes, fonctionnelles d'action et indice des syst\`emes hamiltoniens'', \emph{Bull. Soc. Math. Fr.} \textbf{115} (1987), no. 3, 361--390.
\bibitem[Bru91]{Bru1991} M. Brunella, ``On a theorem of Sikorav'', \emph{Enseign. Math.} \textbf{37} (1991), no. 1-2, 83--87.
\bibitem[Sik87]{Sik1987} J.-C. Sikorav, ``Probl\`emes d'intersections et de points fixes en g\'eom\'etrie hamiltonienne'', \emph{Comment. Math. Helv.} \textbf{62} (1987), no. 1, 62--73.
\end{thebibliography}
\end{document}
